\chapter*{PHỤ LỤC. MÃ NGUỒN VÀ CÁCH TÁI LẬP}
\addcontentsline{toc}{chapter}{PHỤ LỤC. MÃ NGUỒN VÀ CÁCH TÁI LẬP}

\section*{Kho mã nguồn}

Toàn bộ mã nguồn và tài liệu nằm trong kho mã nguồn của nhóm~\cite{nhom13MVTecAnomalyDetection2026}:
\begin{center}
  \url{https://github.com/AnDpTri/mvtec-anomaly-detection}
\end{center}
Kết quả trong Chương 3 được tạo từ commit \texttt{d0216d1}. Sau đó mã nguồn được tổ chức lại mà không đổi thuật toán,
siêu tham số hay cách tính chỉ số; điểm của các mô hình đã lưu, chấm lại trên 48 lượt mẫu, trùng với trước khi tổ chức
lại (sai khác không quá $10^{-7}$). Bảng dưới đây mô tả cấu trúc hiện tại.

\begin{table}[htbp]
  \centering
  \small
  \begin{tabular}{@{}>{\raggedright\arraybackslash}p{6.1cm} >{\raggedright\arraybackslash}p{8.3cm}@{}}
    \toprule
    Thư mục, tệp & Nội dung \\
    \midrule
    \texttt{src/tn/data.py} & Đọc ảnh, mặt nạ và nhãn của MVTec AD, MVTec AD 2, VisA \\
    \texttt{src/tn/download.py} & Tải và giải nén dữ liệu trên Colab \\
    \texttt{src/tn/splits.py} & Tách ảnh tốt thành phần fit và phần validation \\
    \texttt{src/tn/methods/} & Autoencoder, PaDiM, PatchCore, mỗi phương pháp một tệp cùng giao diện: xây dựng mô hình, tính bản đồ bất thường, lưu và nạp mô hình \\
    \texttt{src/tn/metrics.py} & AUROC, AUPRO, ngưỡng và các chỉ số tại ngưỡng \\
    \texttt{tools/build\_tn\_models.py} & Xây dựng, đánh giá và lưu mô hình cho 9 loại \\
    \texttt{tools/make\_colab\_notebook.py} & Sinh notebook Colab có nhúng sẵn mã nguồn \\
    \texttt{tools/tn\_report\_figures.py} & Đo thời gian suy luận, vẽ hình minh họa của Chương 3 \\
    \texttt{src/demo/} & Ứng dụng minh họa Gradio (\texttt{app.py}) và phần nạp mô hình để suy luận (\texttt{inference.py}) \\
    \bottomrule
  \end{tabular}
\end{table}

\section*{Dữ liệu}

MVTec AD (grid, screw, transistor) tải từ bản sao trên Hugging Face bằng lệnh
\begin{verbatim}
python tools/download_mvtec.py <thư mục>/MVTecAD grid screw transistor
\end{verbatim}
MVTec AD 2 (vial, sheet\_metal) tải từ trang
của MVTec, VisA tải từ kho của Amazon; đường tải cụ thể nằm trong \texttt{src/tn/download.py}. Thư mục dữ liệu cần có
\texttt{MVTecAD/}, \texttt{MVTecAD2/} và \texttt{VisA/}.

\section*{Chạy lại thực nghiệm}

Trên Google Colab: sinh notebook bằng \texttt{python tools/make\_colab\_notebook.py}, tải tệp
\texttt{goi\_colab/build\_tn\_models\_colab.ipynb} lên Colab, chọn GPU T4 rồi chạy toàn bộ. Notebook tự tải dữ liệu,
xây dựng mô hình cho 9 loại và ghi kết quả lên Google Drive.

Trên máy có GPU, từ thư mục gốc của kho mã nguồn, sau khi cài các thư viện trong \texttt{requirements.txt}:
\begin{verbatim}
python tools/build_tn_models.py --data <thư mục dữ liệu>
python -m src.demo.app
\end{verbatim}
Lệnh đầu lưu mô hình, ngưỡng, chỉ số và điểm của từng ảnh test, mỗi loại một thư mục con trong
\path{results/checkpoints_tn}, và bỏ qua
các loại đã có kết quả, nên có thể chạy tiếp sau khi bị gián đoạn. Lệnh thứ hai mở ứng dụng minh họa tại
\texttt{http://127.0.0.1:7864}.
